\documentclass[tikz]{standalone}
\usepackage{amsmath}
\usepackage{times}
\usepackage{txfonts}
\usepackage{pgfplots}
\usepackage{csvsimple}
\usetikzlibrary{arrows,intersections,math,calc}
\usetikzlibrary{
    calc,
    arrows.meta,
    decorations.markings,
    decorations.fractals,
    backgrounds,
    positioning,
    shapes.misc
}

\definecolor{jcInnen}{HTML}{FBE9A0} % gelb = Inneres
\definecolor{jcAussen}{HTML}{C9DDF0} % hellblau = Aeusseres
\definecolor{jcKurve}{HTML}{1F3A5F} % Kurvenfarbe
\definecolor{jcAkzent}{HTML}{C0392B} % Hilfslinien
\definecolor{jcGrau}{HTML}{6B6B6B} % Hintergrundgeometrie

\tikzset{
    kurve/.style = {draw=jcKurve, line width=1.1pt},
    kurvedick/.style = {draw=jcKurve, line width=1.4pt},
    strahl/.style = {draw=jcAkzent, line width=0.9pt,
    -{Stealth[length=2.5mm]}},
    tangente/.style = {draw=jcAkzent, line width=0.9pt,
    -{Stealth[length=2mm]}},
    punkt/.style = {circle, fill=jcKurve, inner sep=1.5pt},
    punktrot/.style = {circle, fill=jcAkzent, inner sep=1.5pt},
    schnitt/.style = {cross out, draw=jcAkzent, line width=1pt,
    minimum size=5pt, inner sep=0pt},
    gebietslabel/.style = {font=\large\bfseries, text=jcKurve},
    dashedhelp/.style = {draw=jcGrau, dashed, line width=0.7pt},
    bildtext/.style = {text=jcKurve}
}
\begin{document}
\def\skala{1}
\begin{tikzpicture}[>=latex,thick,scale=\skala, transform shape]
    \fill[jcAussen] (-4.1,-4.3) rectangle (4.1,4.1);
    \fill[jcInnen, smooth, samples=300, domain=0:360, variable=\x]
    plot ({(2.3+0.35*sin(3*\x))*cos(\x)}, {(2.3+0.35*sin(3*\x))*sin(\x)}) -- cycle;
    \draw[line width=22pt, jcAkzent, opacity=0.16, smooth, samples=300,
    domain=0:360, variable=\x]
    plot ({(2.3+0.35*sin(3*\x))*cos(\x)}, {(2.3+0.35*sin(3*\x))*sin(\x)}) -- cycle;
    \foreach \t in {0,15,...,345}{
        \pgfmathsetmacro{\rr}{2.3+0.35*sin(3*\t)}
        \pgfmathsetmacro{\rs}{1.05*cos(3*\t)}
        \pgfmathsetmacro{\tx}{\rs*cos(\t)-\rr*sin(\t)}
        \pgfmathsetmacro{\ty}{\rs*sin(\t)+\rr*cos(\t)}
        \pgfmathsetmacro{\lng}{sqrt(\tx*\tx+\ty*\ty)}
        \pgfmathsetmacro{\nx}{ \ty/\lng}
        \pgfmathsetmacro{\ny}{-\tx/\lng}
        \pgfmathsetmacro{\px}{\rr*cos(\t)}
        \pgfmathsetmacro{\py}{\rr*sin(\t)}
        \draw[jcGrau, line width=0.5pt]
        ({\px-0.38*\nx},{\py-0.38*\ny}) -- ({\px+0.38*\nx},{\py+0.38*\ny});
    }
    \pgfmathsetmacro{\t}{75}
    \pgfmathsetmacro{\rr}{2.3+0.35*sin(3*\t)}
    \pgfmathsetmacro{\rs}{1.05*cos(3*\t)}
    \pgfmathsetmacro{\tx}{\rs*cos(\t)-\rr*sin(\t)}
    \pgfmathsetmacro{\ty}{\rs*sin(\t)+\rr*cos(\t)}
    \pgfmathsetmacro{\lng}{sqrt(\tx*\tx+\ty*\ty)}
    \pgfmathsetmacro{\nx}{ \ty/\lng}
    \pgfmathsetmacro{\ny}{-\tx/\lng}
    \pgfmathsetmacro{\px}{\rr*cos(\t)}
    \pgfmathsetmacro{\py}{\rr*sin(\t)}
    \draw[tangente, line width=1.1pt]
    ({\px},{\py}) -- ({\px+0.55*\nx},{\py+0.55*\ny});
    \node[text=jcAkzent]
    at ({\px+1.05*\nx},{\py+1.05*\ny}) {$n(t)$};
    \draw[kurvedick, smooth, samples=300, domain=0:360, variable=\x]
    plot ({(2.3+0.35*sin(3*\x))*cos(\x)}, {(2.3+0.35*sin(3*\x))*sin(\x)}) -- cycle;
    \node[gebietslabel] at (0,0) {$I$};
    \node[gebietslabel] at (3.4,3.4) {$A$};
    \node[text=jcAkzent] at (0,-3.6)
    {$T_\varepsilon=\varphi\bigl(S^1\times(-\varepsilon,\varepsilon)\bigr)$};
\end{tikzpicture}
\end{document}

